\documentclass[11pt]{article}

\usepackage[margin=1in]{geometry}
\usepackage{amsmath,amssymb,amsthm}
\usepackage{mathtools}
\usepackage{bm}
\usepackage{enumitem}
\usepackage{hyperref}
\usepackage{xcolor}

\newcommand{\Z}{\mathcal{Z}}
\newcommand{\X}{\mathcal{X}}
\newcommand{\Obs}{\mathcal{O}}
\newcommand{\GFI}{\textsc{gfi}}
\newcommand{\copyact}{\textsf{copy}}
\newcommand{\subact}{\textsf{sub}}
\newcommand{\inssact}{\textsf{ins}}
\newcommand{\delact}{\textsf{del}}
\newcommand{\given}{\mid}
\newcommand{\pLM}{p_{\mathrm{LM}}}
\newcommand{\dDL}{d_{\mathrm{DL}}}
\newcommand{\lse}{\operatorname{logsumexp}}
\newcommand{\Cat}{\mathrm{Cat}}
\newcommand{\Dir}{\mathrm{Dir}}
\newcommand{\eos}{\textsc{eos}}
\newcommand{\Vocab}{\mathbb{V}}
\newcommand{\Real}{\mathbb{R}}
\newcommand{\Ind}{\mathbb{1}}
\newcommand{\code}[1]{\texttt{#1}}

\newtheorem{theorem}{Theorem}
\newtheorem{lemma}{Lemma}
\newtheorem{proposition}{Proposition}
\newtheorem{definition}{Definition}
\newtheorem{remark}{Remark}

\title{A Token-Level Noisy-Channel Model of Language Comprehension:\\
Generative Semantics, Sequential Monte Carlo, and Rejuvenation}
\author{genjax noisy-channel port \\ \small \texttt{src/genjax\_port/}}
\date{\today}

\begin{document}
\maketitle

\begin{abstract}
We give a formal description of the genjax-native noisy-channel model implemented in
\code{src/genjax\_port/}. A speaker draws an \emph{intended} token sequence from a language-model
prior and a word-level noisy channel corrupts it into the \emph{observed} sentence; inference is
the posterior over intended sequences. We state the model as a genjax generative function, define
the inference problem in the conventions of the genjax Generative Function Interface (\GFI), and
then prove the two correctness claims that the implementation relies on: (i) the word-scan
sequential Monte Carlo (SMC) filter is \emph{properly weighted} for the sequence of posteriors and
its normalizing-constant estimator is unbiased, with the per-word incremental weight collapsing to
a $\lse$ of branch evidences; and (ii) the rejuvenation (incremental-reanalysis) moves built on
genjax's \code{Rejuvenate} edit request are valid SMCP3 / Metropolis--Hastings moves that leave the
target posterior invariant. Notation follows the genjax library where applicable; the appendix
recalls the \GFI{} weight identities used in the proofs.
\end{abstract}

\tableofcontents

\section{Notation and the generative function interface}
\label{sec:gfi}

We use the conventions of the genjax library. A \emph{generative function} $p$ over arguments
$a$ defines a probability measure on a space of \emph{choice maps} $t$ (finite maps from
\emph{addresses} to values) together with a deterministic return value. Writing $p(t; a)$ for its
density with respect to the appropriate base measure, the \GFI{} methods used below have the
following semantics (see Appendix~\ref{app:gfi} and \cite{genjax,gen}).

\begin{itemize}[leftmargin=1.4em]
  \item \code{simulate}$(k, a) \to (t, \cdot)$ samples $t \sim p(\cdot; a)$ and records the
        \emph{score} $\log p(t; a)$.
  \item \code{generate}$(k, c, a) \to (t, w)$ (called \code{importance} for distributions) samples
        the unconstrained addresses from an internal proposal $q$ given the constraint choice map
        $c \subseteq t$, and returns the \emph{importance weight}
        \begin{equation}\label{eq:gen-weight}
          w \;=\; \log \frac{p(t; a)}{q(t_{\setminus c} \given c; a)},
        \end{equation}
        where $t_{\setminus c}$ are the addresses sampled by $q$. When $c$ constrains \emph{all}
        random choices, $q$ is empty and $w = \log p(t;a) = \log p(c; a)$.
  \item \code{assess}$(t, a) \to (\log p(t; a), \mathrm{ret})$ scores a fully-constrained $t$.
  \item \code{edit}$(k, t, r, \delta a) \to (t', w, \delta\mathrm{ret}, r')$ applies an edit request
        $r$ (e.g.\ \code{Update}, \code{Rejuvenate}) and argument diffs $\delta a$, returning a new
        trace $t'$, an \emph{edit weight} $w$, and a \emph{backward request} $r'$ that reverses the
        move. For an \code{Update} that overwrites the addresses in a constraint $u$ (with
        $\delta a$ a no-change),
        \begin{equation}\label{eq:update-weight}
          w \;=\; \log \frac{p(t'; a)}{p(t; a)},
          \qquad r' = \code{Update}(t|_{\mathrm{dom}(u)}),
        \end{equation}
        i.e.\ the backward request restores the overwritten addresses to their old values.
\end{itemize}

These identities are all we need; they are restated with derivations in Appendix~\ref{app:gfi}.

\section{The model}
\label{sec:model}

\subsection{Observed data and deterministic segmentation}

Let $\Vocab$ be the BPE vocabulary ($|\Vocab| = V$) of a fixed autoregressive language model with
next-token distribution $\pLM(\cdot \given u)$ for a context $u \in \Vocab^{*}$. The observation is a
token sequence which a deterministic, data-only segmenter (\code{noise\_word.segment\_words})
splits into $W$ \emph{words}
\begin{equation}
  o \;=\; (w_1, \dots, w_W), \qquad w_i = (o_{i,1}, \dots, o_{i,n_i}) \in \Vocab^{n_i}.
\end{equation}
Because segmentation is a deterministic function of $o$, every particle consumes the same word
boundaries; the words are \emph{indexing data}, not latent variables.

\subsection{Latent variables}

The latent state is the \emph{intended} token sequence together with a per-word alignment. For each
word $i$ we introduce an \emph{action} $a_i$ and the intended tokens it emits. In the
substitution-only model (the implemented M1 filter, \code{smc\_substitution.py}) the action set is
\begin{equation}
  \mathcal{A}_i \;=\; \{\copyact\} \,\cup\, \{\subact_x : x \in \mathcal{S}_i\},
\end{equation}
where $\mathcal{S}_i \subseteq \Vocab$ is the (data-derived, finite) set of single-token
substitution candidates for $w_i$ returned by \code{word\_sub\_candidates}: word-initial single
vocabulary tokens within Damerau--Levenshtein distance $\le D$ of the surface string of $w_i$.
Section~\ref{sec:extensions} adds insertion ($\inssact$) and deletion ($\delact$).

Given an action, the emitted intended tokens are deterministic:
\begin{equation}
  \mathrm{emit}(a_i) =
  \begin{cases}
    (o_{i,1}, \dots, o_{i,n_i}) & a_i = \copyact \quad (\text{the literal word, $n_i$ tokens}),\\[2pt]
    (x) & a_i = \subact_x \quad (\text{one token}).
  \end{cases}
\end{equation}
The intended sequence is the concatenation $z = \mathrm{emit}(a_1) \,\Vert\, \cdots \,\Vert\,
\mathrm{emit}(a_W)$, with a fixed start-of-sequence seed $z_0 = \eos$. We write $z_{<i}$ for the
prefix emitted by words $1,\dots,i-1$ (the LM context entering word $i$).

\subsection{Priors and channel}

\paragraph{Action prior.} A global mixing weight $\pi \in \Delta^{2}$ (over the abstract action
\emph{kinds} $\copyact,\subact,\inssact$) is drawn $\pi \sim \Dir(\alpha)$ with
$\alpha = \code{ACTION\_ALPHAS} = (3,1,1)$, and each action kind is scored under $\pi$. We write
$\kappa(a_i) \in \{\copyact,\subact,\inssact\}$ for the kind of $a_i$.

\paragraph{Language-model prior.} The emitted intended tokens are scored autoregressively under the
LM. For word $i$ contributing tokens $\mathrm{emit}(a_i) = (y_1,\dots,y_m)$ after context $z_{<i}$,
\begin{equation}\label{eq:lm-word}
  \log \pLM\!\big(\mathrm{emit}(a_i) \given z_{<i}\big)
  \;=\; \sum_{r=1}^{m} \log \pLM\!\big(y_r \given z_{<i} \Vert y_{1:r-1}\big).
\end{equation}

\paragraph{Channel.} The word-level channel likelihood $P(w_i \given a_i)$ is
\begin{equation}\label{eq:channel}
  \log P(w_i \given a_i) =
  \begin{cases}
    0 & a_i = \copyact \quad(\text{exact match}),\\[2pt]
    \dDL\!\big(\mathrm{surf}(x), \mathrm{surf}(w_i)\big)\cdot \log \theta & a_i = \subact_x,
  \end{cases}
\end{equation}
with substitution sharpness $\theta = \code{SUB\_PARAM} = 0.1$ and $\dDL$ the
Damerau--Levenshtein character distance (\code{noise\_word.word\_sub\_loglik}).

\subsection{Joint density}

Collecting the above, the joint density of the global weight $\pi$, the actions $a_{1:W}$, and the
data $o$ (with $z$ and the word boundaries deterministic given $a_{1:W}$ and $o$) is
\begin{equation}\label{eq:joint}
  p(\pi, a_{1:W}, o)
  \;=\; \Dir(\pi; \alpha)\;
  \prod_{i=1}^{W}
  \underbrace{\pi_{\kappa(a_i)}\;
  \pLM\!\big(\mathrm{emit}(a_i)\given z_{<i}\big)\;
  P(w_i \given a_i)}_{\textstyle =:\,\exp E_i(a_i)} .
\end{equation}
We call $E_i(c)$ the \emph{branch evidence} of candidate $c \in \mathcal{A}_i$ at word $i$; it is the
log-density contribution of explaining $w_i$ by $c$ given the particle's prefix:
\begin{equation}\label{eq:branch-ev}
  E_i(c) \;=\; \log \pi_{\kappa(c)}
  \;+\; \log \pLM\!\big(\mathrm{emit}(c)\given z_{<i}\big)
  \;+\; \log P(w_i \given c).
\end{equation}

\subsection{Realization as a genjax generative function}
\label{sec:genfn}

The model is the word scan whose per-step kernel is the genjax \code{Switch} combinator
$\code{make\_word\_model}(n_i, |\mathcal{S}_i|)$ over the branches $\mathcal{A}_i$:
branch $0$ ($\copyact$) samples $n_i$ language-model tokens at addresses
$\code{t0},\dots,\code{t}(n_i{-}1)$; branch $k\ge 1$ ($\subact_{x_k}$) samples one token at
$\code{t0}$; every branch closes with a deterministic factor address $\code{ch}$ carrying the
action-prior$+$channel term. Concretely, each branch is a generative function whose \code{assess}
score (Appendix~\ref{app:gfi}) equals \eqref{eq:branch-ev}; this is verified numerically against the
manual right-hand side in \code{tests/test\_word\_model.py} (branch \code{importance}
$=$ manual joint, to $10^{-2}$). The token choices use the custom \code{exact\_density} distribution
\code{lm\_token} that wraps the in-graph LM (\code{lm\_genjax.py}). The branch index is an
addressed categorical choice; \emph{this is the address that rejuvenation edits}
(Section~\ref{sec:rejuv}).

\begin{remark}[Single-token substitution / deletion]
Both $\subact$ and (later) $\delact$ are $N{:}1$: the intended word is a single BPE token. This is
the implemented scope; multi-token intended words ($M{:}N$) are future work and do not affect any
statement below.
\end{remark}

\section{Inference target and the SMC filter}
\label{sec:smc}

\subsection{Sequence of targets}

Let $\xi_{1:i} := (\pi, a_{1:i})$ denote the latents instantiated after the first $i$ words. Define
the \emph{unnormalized} step-$i$ target as the joint density of the prefix and the first $i$ observed
words,
\begin{equation}\label{eq:gamma}
  \gamma_i(\xi_{1:i}) \;=\; \Dir(\pi;\alpha)\prod_{j=1}^{i}\exp E_j(a_j),
  \qquad
  Z_i \;=\; \sum_{a_{1:i}}\int \gamma_i \,d\pi \;=\; p(w_{1:i}),
\end{equation}
with $\gamma_0 = \Dir(\pi;\alpha)$ and $Z_0 = 1$. The normalized target is
$\bar\gamma_i = \gamma_i / Z_i = p(\xi_{1:i}\given w_{1:i})$, and the final target $\bar\gamma_W$ is the
posterior of interest. The genjax counterpart is a \code{Target} whose generative function is the
word scan and whose constraints fix the observed tokens; the implementation realizes the sequence
$\bar\gamma_0,\dots,\bar\gamma_W$ by extending and reweighting particles word by word.

\subsection{Proposal and incremental weight}

The filter uses the \emph{locally optimal} (``data-driven'') proposal: holding $\pi$ fixed at its
prior draw,\footnote{The implemented filter samples $\pi$ once from the prior per particle and does
not update it during the forward pass; this is a deliberate approximation (a parameter rejuvenation
move is the optional R4 extension). All SMC statements below hold for the model with $\pi$ treated as
part of $\xi_1$ and proposed from the prior at step~1; subsequent steps leave $\pi$ unchanged.}
each particle extends its alignment at word $i$ by sampling a branch $c \in \mathcal{A}_i$ from
\begin{equation}\label{eq:proposal}
  q_i(c \given \xi_{1:i-1}) \;=\; \frac{\exp E_i(c)}{\sum_{c'\in\mathcal{A}_i}\exp E_i(c')}
  \;=\; \mathrm{softmax}_c\, E_i(c),
\end{equation}
implemented by \code{proposal.propose} on the evidence vector
$E_i(\cdot)$ assembled by \code{smc\_substitution.word\_log\_evidence} (which equals the per-branch
\code{Switch} \code{importance} weights, \S\ref{sec:genfn}). The corresponding incremental importance
weight is the ratio of the target increment to the proposal:
\begin{align}
  w_i
  &\;=\; \frac{\gamma_i(\xi_{1:i})}{\gamma_{i-1}(\xi_{1:i-1})\,\cdot\,q_i(a_i\given \xi_{1:i-1})}
  \;=\; \frac{\exp E_i(a_i)}{q_i(a_i\given\xi_{1:i-1})}. \label{eq:wstep}
\end{align}

\begin{lemma}[Weight collapse]\label{lem:collapse}
With the locally optimal proposal \eqref{eq:proposal}, the incremental weight \eqref{eq:wstep} is
independent of the sampled branch $a_i$ and equals the per-word evidence sum:
\begin{equation}
  \log w_i \;=\; \lse_{c\in\mathcal{A}_i}\, E_i(c)
  \;=\; \log\!\!\sum_{c\in\mathcal{A}_i}\exp E_i(c).
\end{equation}
\end{lemma}
\begin{proof}
Substitute \eqref{eq:proposal} into \eqref{eq:wstep}:
$
  w_i = \exp E_i(a_i) \big/ \big(\exp E_i(a_i)/\textstyle\sum_{c}\exp E_i(c)\big)
      = \sum_{c}\exp E_i(c).
$
The $\exp E_i(a_i)$ factors cancel, leaving a quantity that does not depend on $a_i$. Taking logs
gives the $\lse$. \qed
\end{proof}

This is exactly the line
\code{log\_marginal += logsumexp(log\_w) - log(P)} together with
\code{log\_w = logsumexp(log\_ev)} in \code{smc\_substitution.py} (and \code{proposal.py}); note
$\log w_i$ is identical across particles \emph{up to the prefix dependence} of $E_i$, so different
particles still receive different weights through their distinct contexts $z_{<i}$.

\subsection{Resampling and the algorithm}

After computing $\{w_i^{(p)}\}_{p=1}^{P}$ the filter resamples ancestors
$A^{(p)} \sim \Cat\big(\mathrm{softmax}_p \log w_i^{(p)}\big)$ (multinomial; \code{jax.random.categorical}
on normalized log-weights) and copies the corresponding particle state. Algorithm~\ref{alg:smc}
summarizes one sweep.

\begin{figure}[t]
\hrule\vspace{2pt}
\textbf{Algorithm 1: word-scan SMC (substitution model).}
\begin{enumerate}[leftmargin=2em,itemsep=1pt]
  \item Draw $\pi^{(p)}\sim\Dir(\alpha)$; set $z^{(p)}_0=\eos$, $\Ind^{(p)}=1$, $\widehat Z = 1$.
  \item For $i = 1,\dots,W$:
  \begin{enumerate}[leftmargin=1.4em,itemsep=0pt]
    \item For each particle, form $E_i(\cdot)$ from \eqref{eq:branch-ev} (one LM forward $+$ gather;
          \code{word\_log\_evidence}).
    \item Sample $a_i^{(p)} \sim q_i(\cdot)$ \eqref{eq:proposal} and set
          $\log w_i^{(p)} = \lse_c E_i^{(p)}(c)$ (Lemma~\ref{lem:collapse}).
    \item Update $\widehat Z \mathrel{{*}{=}} \tfrac1P\sum_p w_i^{(p)}$; resample ancestors
          $A^{(p)}\sim\Cat(\mathrm{softmax}\,\log w_i^{(\cdot)})$; copy state; emit
          $\mathrm{emit}(a_i)$ into each $z^{(p)}$.
  \end{enumerate}
  \item Return the particles $\{z^{(p)}\}$ (posterior sample) and $\log\widehat Z$.
\end{enumerate}
\hrule
\label{alg:smc}
\end{figure}

\subsection{Correctness}

\begin{definition}[Proper weighting]
A weighted particle set $\{(\xi^{(p)}, w^{(p)})\}_{p=1}^P$ is \emph{properly weighted} for an
unnormalized target $\gamma$ with constant $Z=\int\gamma$ if for all integrable $f$,
$\;\mathbb{E}\!\big[\tfrac1P\sum_p w^{(p)} f(\xi^{(p)})\big] = \int f\,\gamma = Z\,\mathbb{E}_{\bar\gamma}[f].$
\end{definition}

\begin{theorem}[Properly weighted filter and unbiased evidence]\label{thm:smc}
For every $i$, the weighted particles produced after step $i$ of Algorithm~\ref{alg:smc} are properly
weighted for $\gamma_i$ in \eqref{eq:gamma}. Consequently the after-resampling particles are
distributed (in the proper-weighting sense) according to the posterior $\bar\gamma_i$, and
\begin{equation}
  \mathbb{E}\big[\widehat Z_i\big] \;=\; Z_i \;=\; p(w_{1:i}),
  \qquad
  \widehat Z_i \;=\; \prod_{j=1}^{i}\Big(\tfrac1P\textstyle\sum_p w_j^{(p)}\Big).
\end{equation}
In particular $\widehat Z_W$ is an unbiased estimator of the marginal likelihood $p(o)$, matching the
genjax \code{ParticleCollection} estimator $\log\widehat Z = \sum_i\big(\lse_p \log w_i^{(p)} - \log P\big)$
\textup{(\code{smc.py})}.
\end{theorem}

\begin{proof}
Induction on $i$. \emph{Base.} At $i=0$ each particle has weight $1$ and $\pi^{(p)}\sim\Dir(\alpha)=\bar\gamma_0$, which is properly weighted for $\gamma_0$ with $Z_0=1$.

\emph{Step.} Assume proper weighting for $\gamma_{i-1}$. The extension samples $a_i \sim q_i(\cdot\given\xi_{1:i-1})$ and multiplies the weight by $w_i$ of \eqref{eq:wstep}. This is one step of sequential importance sampling with proposal $q_i$ and target increment $\gamma_i/\gamma_{i-1} = \exp E_i(a_i)$; the standard SIS identity gives, for any $f$,
\begin{align*}
  \mathbb{E}\Big[\tfrac1P\!\sum_p w_{i-1}^{(p)} w_i^{(p)} f(\xi_{1:i}^{(p)})\Big]
  &= \mathbb{E}\Big[\tfrac1P\!\sum_p w_{i-1}^{(p)} \,\mathbb{E}_{a_i\sim q_i}\!\Big[\tfrac{\exp E_i(a_i)}{q_i(a_i)} f\Big]\Big]\\
  &= \mathbb{E}\Big[\tfrac1P\!\sum_p w_{i-1}^{(p)} \!\!\sum_{a_i}\!\exp E_i(a_i)\, f(\xi_{1:i}^{(p)})\Big]
   = \textstyle\int f\,\gamma_i,
\end{align*}
where the last equality uses $q_i(a_i)>0$ on $\mathrm{supp}\,\exp E_i$ (every branch with positive
evidence has positive proposal mass, since $q_i\propto \exp E_i$) and the inductive hypothesis with
$\gamma_i = \gamma_{i-1}\exp E_i$. Hence the extended set is properly weighted for $\gamma_i$.

\emph{Resampling.} Multinomial resampling with probabilities $\propto w_i^{(p)}$ replaces weights by
the mean $\bar w_i = \tfrac1P\sum_p w_i^{(p)}$ and is unbiased: $\mathbb{E}[f \text{ after}] =
\mathbb{E}[\tfrac1P\sum_p \tfrac{w_i^{(p)}}{\bar w_i} f]\cdot \bar w_i$ in expectation preserves
$\int f\,\gamma_i$ (a standard SMC result, e.g.\ \cite{chopin2020}). Folding $\bar w_i$ into
$\widehat Z_i = \widehat Z_{i-1}\,\bar w_i$ keeps the product unbiased for $Z_i$ by the tower rule.
The telescoping product is $\widehat Z_i = \prod_{j\le i}\bar w_j$, and by Lemma~\ref{lem:collapse}
$\log \bar w_j = \lse_p \log w_j^{(p)} - \log P$. \qed
\end{proof}

\begin{remark}[Why edits, not re-derivations]
Theorem~\ref{thm:smc} only requires (a) $q_i\propto\exp E_i$ on the support and (b) the weight
\eqref{eq:wstep}. The genjax realization obtains both from the \GFI: $E_i(c)$ is the \code{assess}
score of branch $c$, and \eqref{eq:gen-weight} with the fully-constrained branch gives exactly
$w$ for the chosen branch; the $\lse$ collapse is the proposal normalizer. Thus the hand-rolled
outer loop and the native model agree by construction, not by coincidence.
\end{remark}

\section{Rejuvenation: incremental reanalysis}
\label{sec:rejuv}

Resampling repeatedly prunes alternatives, so an early word committed under local information can
never be revised once later context disambiguates it. Rejuvenation interleaves, after each
resampling step, MCMC moves that revisit earlier words within a \emph{lookback} window. Crucially,
these moves target the \emph{same} posterior $\bar\gamma_i$; they do not change the target, they let
the chain reach high-posterior reanalyses and re-diversify the particle set.

\subsection{The \texttt{Rejuvenate} edit request (SMCP3)}

genjax's \code{Rejuvenate} (\code{inference/requests/rejuvenate.py}) implements an SMCP3 move
\cite{smcp3} from a single proposal $Q$ used for both the forward ($K$) and backward ($L$) kernels.
Let $t$ be a particle's trace, $\mathcal{B}\subseteq\mathrm{dom}(t)$ the edited addresses (e.g.\ the
action and token addresses aligned to a revisited word), and $\phi(\cdot)$ the
\code{argument\_mapping} producing $Q$'s arguments from a choice map. The move is:
\begin{enumerate}[leftmargin=1.8em,itemsep=1pt]
  \item \textbf{Forward.} Draw $\zeta \sim Q(\cdot;\phi(t))$ with score $s_{\mathrm{fwd}} = \log Q(\zeta;\phi(t))$.
  \item \textbf{Update.} Apply $\code{Update}(\zeta)$: $t' = t$ with $\mathcal{B}$ overwritten by
        $\zeta$, with edit weight (cf.\ \eqref{eq:update-weight}, $\delta a$ a no-change)
        \begin{equation}\label{eq:wupd}
          w_{\mathrm{upd}} = \log\frac{p(t';a)}{p(t;a)},
        \end{equation}
        and backward request $\code{Update}(t|_{\mathcal B})$ restoring the old values.
  \item \textbf{Backward.} Score the old values under $Q$ at the \emph{new} point:
        $s_{\mathrm{bwd}} = \log Q\big(t|_{\mathcal B};\,\phi(t')\big)$.
  \item \textbf{Weight.} Return the SMCP3 weight
        \begin{equation}\label{eq:smcp3}
          \mathcal{W} \;=\; w_{\mathrm{upd}} + s_{\mathrm{bwd}} - s_{\mathrm{fwd}}
          \;=\; \log\frac{p(t';a)\,Q\big(t|_{\mathcal B};\phi(t')\big)}{p(t;a)\,Q\big(\zeta;\phi(t)\big)}.
        \end{equation}
\end{enumerate}

There are two valid ways to consume $\mathcal{W}$, both implemented/available:

\paragraph{(A) Metropolis--Hastings.} Accept $t'$ with probability $\min(1,e^{\mathcal W})$ (i.e.\
accept iff $\log u < \mathcal W$, $u\sim\mathrm{Unif}(0,1)$), else keep $t$. Particle weights are
untouched. This is the default in the spikes and tests
(\code{tests/test\_noisy\_channel.py}: $\;\code{accept = log(uniform) < w}$).

\paragraph{(B) SMCP3 reweighting.} Move deterministically to $t'$ and multiply the particle weight by
$e^{\mathcal W}$ (no accept/reject). This keeps the collection properly weighted (Theorem~\ref{thm:rejuv-smcp3}).

\subsection{Correctness of the MH consumption}

\begin{theorem}[MH invariance]\label{thm:mh}
Let $\Pi(t \to t')$ be the kernel that proposes $\zeta\sim Q(\cdot;\phi(t))$, forms $t'$ by
overwriting $\mathcal B$ with $\zeta$, and accepts with probability $\min(1,e^{\mathcal W})$ for
$\mathcal W$ in \eqref{eq:smcp3}. Then $\Pi$ satisfies detailed balance with respect to the posterior
$\bar\gamma \propto p(\cdot\given o)$, hence leaves $\bar\gamma$ invariant.
\end{theorem}

\begin{proof}
Because the move overwrites $\mathcal B$ and leaves $\mathrm{dom}(t)\setminus\mathcal B$ fixed, the
proposal of $t'$ from $t$ has density $Q(\zeta;\phi(t))$ with $\zeta = t'|_{\mathcal B}$, and the
reverse proposal of $t$ from $t'$ has density $Q(t|_{\mathcal B};\phi(t'))$ — exactly the $L$-kernel
genjax scores as $s_{\mathrm{bwd}}$. The posterior density satisfies
$\bar\gamma(t)\propto p(t;a)$ (the $o$-addresses are constrained and common to $t,t'$, so they cancel
in any ratio). The Metropolis--Hastings acceptance for proposal $Q$ and target $\bar\gamma$ is
\begin{equation}
  \alpha(t\to t') = \min\!\Big(1,\;\frac{\bar\gamma(t')\,Q(t|_{\mathcal B};\phi(t'))}{\bar\gamma(t)\,Q(t'|_{\mathcal B};\phi(t))}\Big)
  = \min\!\big(1, e^{\mathcal W}\big),
\end{equation}
since $\bar\gamma(t')/\bar\gamma(t) = p(t';a)/p(t;a) = e^{w_{\mathrm{upd}}}$ by \eqref{eq:wupd}, and
the $Q$-ratio is $e^{s_{\mathrm{bwd}}-s_{\mathrm{fwd}}}$. Detailed balance
$\bar\gamma(t)\,Q(t'|_{\mathcal B};\phi(t))\,\alpha(t\to t') = \bar\gamma(t')\,Q(t|_{\mathcal B};\phi(t'))\,\alpha(t'\to t)$
is the standard Metropolis--Hastings identity, which holds verbatim once the proposal and target
ratios are identified as above. Invariance follows by summing/integrating detailed balance over
$t$. \qed
\end{proof}

\begin{remark}[Why $\mathcal W$ is exactly the MH ratio]
The content of Theorem~\ref{thm:mh} is that genjax assembles the three \GFI{} quantities
$w_{\mathrm{upd}}, s_{\mathrm{fwd}}, s_{\mathrm{bwd}}$ so that \eqref{eq:smcp3} \emph{is} the MH
log-acceptance ratio. The \code{Update} weight \eqref{eq:update-weight} supplies the target ratio and
the backward request supplies precisely the addresses $L$ must score; no separate hand-derivation of
the ratio is needed (contrast the by-hand involutive-MH route in \code{src/gen\_inference.jl}).
\end{remark}

\subsection{Correctness of the SMCP3 (reweighting) consumption}

\begin{theorem}[SMCP3 proper weighting]\label{thm:rejuv-smcp3}
Suppose $\{(t^{(p)}, w^{(p)})\}$ is properly weighted for the posterior $\bar\gamma$ (unnormalized
$p(\cdot;a)$ on the latent addresses, $o$ constrained). Applying the deterministic update
$t^{(p)}\!\mapsto t'^{(p)}$ and reweighting $w^{(p)}\!\mapsto w^{(p)}e^{\mathcal W^{(p)}}$ with
$\mathcal W$ from \eqref{eq:smcp3} yields a set properly weighted for the same $\bar\gamma$.
\end{theorem}

\begin{proof}
This is the SMCP3 theorem of \cite{smcp3} specialized to a deterministic forward map with auxiliary
randomness $\zeta\sim Q(\cdot;\phi(t))$ and backward auxiliary scored by $Q(\cdot;\phi(t'))$. For any
$f$,
\begin{align*}
  \mathbb{E}\Big[\tfrac1P\!\sum_p w^{(p)} e^{\mathcal W^{(p)}} f(t'^{(p)})\Big]
  &= \mathbb{E}\Big[\tfrac1P\!\sum_p w^{(p)} \,\mathbb{E}_{\zeta\sim Q(\cdot;\phi(t^{(p)}))}\!\Big[ \tfrac{p(t'^{(p)})\,Q(t^{(p)}|_{\mathcal B};\phi(t'^{(p)}))}{p(t^{(p)})\,Q(\zeta;\phi(t^{(p)}))} f(t'^{(p)})\Big]\Big].
\end{align*}
Using proper weighting to replace the outer $\tfrac1P\sum_p w^{(p)}(\cdot)$ by
$\int(\cdot)\,p(t)\,dt$ and then expanding the inner expectation as an integral over $\zeta$ with
density $Q(\zeta;\phi(t))$, the $p(t)$ and $Q(\zeta;\phi(t))$ factors cancel:
\[
  = \int\!\!\int p(t'(t,\zeta))\,Q(t|_{\mathcal B};\phi(t'))\, f(t'(t,\zeta))\, d\zeta\,dt
  = \int p(t')\, f(t')\, dt' = \textstyle\int f\,\bar\gamma\,,
\]
where the penultimate step is the change of variables $(t,\zeta)\mapsto t'$ whose reverse map is
scored by $Q(\cdot;\phi(t'))$ (the $L$-kernel), with unit Jacobian for the discrete overwrite. Hence
proper weighting for $\bar\gamma$ is preserved. \qed
\end{proof}

\subsection{The move families}
\label{sec:moves}

\paragraph{R1 — substitution flip (fixed dimension).} Edit set $\mathcal B$ = the action address and
its emitted token address for a revisited word $i$. The proposal $Q$ is the local posterior over that
word's candidate set, reweighted by the live LM:
\begin{equation}\label{eq:r1prop}
  Q(c;\phi(t)) \;=\; \mathrm{softmax}_{c\in\mathcal A_i}\, E_i(c),
\end{equation}
identical in form to the forward proposal \eqref{eq:proposal} but evaluated against the
\emph{current} (possibly later-updated) context. Because $\mathcal B$ has fixed cardinality, the
overwrite is dimension-preserving and Theorem~\ref{thm:mh}/\ref{thm:rejuv-smcp3} apply directly. The
spike in \code{tests/test\_noisy\_channel.py} exhibits this move flipping a literal reading
$x=\textsf{`` too''}$ to the posterior-preferred $\textsf{`` to''}$ under \eqref{eq:r1prop}.

\paragraph{R2 — add/delete (trans-dimensional).} A reversible-jump move that inserts an omitted
intended word ($\delact$ reversed) or deletes a posited one, changing the length of $z$. The forward
draws a coin (add vs.\ delete) and, on add, a new intended token from the LM-restricted proposal; the
backward maps add$\leftrightarrow$delete so the move is its own reverse family. The validity is again
\eqref{eq:smcp3}: \code{Update} reports $w_{\mathrm{upd}} = \log p(t')/p(t)$ across the changed
suffix, and the add/delete symmetry makes $\phi$ produce matching $K$/$L$ arguments. As the emitted
tokens are discrete and regenerated (not a continuous reparametrization), the Jacobian is unity and no
extra term enters \eqref{eq:smcp3}; correctness is Theorem~\ref{thm:mh} with $\mathcal B$ the changed
addresses plus the move-type auxiliary.

\paragraph{R3 — surprisal-gated scheduling.} The \emph{choice of whether/where} to rejuvenate (a
surprisal-conditioned trigger over a lookback window, ported from \code{gen\_inference.jl}) is part of
the kernel \emph{schedule}, not the target. Any data-dependent but state-independent schedule of
$\bar\gamma$-invariant moves is itself $\bar\gamma$-invariant; gating only changes mixing, not the
stationary distribution.

\subsection{The cost of reanalysis}

The LM is causal: overwriting an intended token at position $k$ changes
$\pLM(\cdot\given\text{context})$ for all later positions. The target ratio in $w_{\mathrm{upd}}$
therefore re-scores the suffix from $k$ to the current frontier $N$; the prefix cancels. With a
lookback window of size $\ell$ we have $k \ge N-\ell$, so each move costs $O(\ell)$ LM evaluations —
$\ell$ is the reach/cost knob, and re-scoring the suffix is exactly what lets later context vote on
word $k$. This is why the rejuvenation-bearing trace must keep the suffix tokens as \emph{live}
addressed LM choices (Section~\ref{sec:genfn}), not as frozen factors.

\section{Extensions present in the unified filter}
\label{sec:extensions}

The deployed reference filter (\code{particle\_filter\_unified.py}) augments $\mathcal A_i$ with two
further actions, with identical SMC/rejuvenation theory (only $E_i$ and $\mathrm{emit}$ change):
\begin{itemize}[leftmargin=1.4em,itemsep=1pt]
  \item \textbf{Insertion} ($\inssact$): $w_i$ is spurious; $\mathrm{emit}(\inssact)=()$ and
        $\log P(w_i\given\inssact) = n_i\cdot(-\log V)$ (a flat per-token cost). It adds one column to
        $E_i$ with branch evidence $\log\pi_{\inssact} + n_i(-\log V)$.
  \item \textbf{Deletion} ($\delact$): an omitted intended single-token word is posited \emph{before}
        $w_i$, proposed lookahead-guided toward $o_{i,1}$. This is a within-step pre-scan whose
        importance contribution multiplies into $w_i$; with the lookahead proposal the same
        weight-collapse argument (Lemma~\ref{lem:collapse}) applies to the augmented branch set.
\end{itemize}
In the genjax-native port these correspond to additional \code{Switch} branches (insertion: a branch
emitting zero tokens; deletion: a masked variable-length pre-emission), handled by the \code{Mask}
combinator; the proofs above are unchanged because they only use the branch evidences $E_i$ and the
\GFI{} weight identities.

\appendix
\section{Generative function interface weight identities}
\label{app:gfi}

For completeness we restate the \GFI{} weights used above; these are the genjax/Gen semantics
\cite{genjax,gen}. Let $p(t;a)$ be the density of a generative function and let a choice map split as
$t = (c, u)$ into constrained $c$ and unconstrained $u$.

\paragraph{\texttt{generate}/\texttt{importance}.} With internal (or supplied) proposal $q(u\given c;a)$,
\[
  (t,w)\sim\code{generate}(k,c,a),\qquad
  w = \log p(t;a) - \log q(u\given c;a).
\]
If $c=t$ (fully constrained), $w=\log p(t;a)$. For a distribution wrapped by \code{exact\_density}
(our \code{lm\_token}, \code{obs\_dist}, \code{factor}), $p$ is the supplied log-density and the
fully-constrained weight equals \code{assess}.

\paragraph{\texttt{update} edit.} For an \code{Update} request with constraint $u'$ over addresses
$\mathcal B=\mathrm{dom}(u')$ and no-change argdiffs,
\[
  (t',w,\,\cdot\,,r')=\code{edit}(k,t,\code{Update}(u'),\,\delta a{=}\mathrm{nochange}),
  \quad w=\log\frac{p(t';a)}{p(t;a)},\quad r'=\code{Update}(t|_{\mathcal B}).
\]
The backward request $r'$ restores the overwritten addresses; this is what supplies the $L$-kernel's
scoring point in \code{Rejuvenate}. (When $\delta a\ne\mathrm{nochange}$, $w$ additionally accounts
for the argument change; the spikes pass \code{Diff.no\_change(tr.get\_args())} so the simple form
applies.)

\paragraph{\texttt{Rejuvenate} edit.} As derived in Section~\ref{sec:rejuv}, with proposal $Q$ and
\code{argument\_mapping} $\phi$,
\[
  \mathcal W = w_{\mathrm{upd}} + \log Q\big(t|_{\mathcal B};\phi(t')\big) - \log Q\big(\zeta;\phi(t)\big),
\]
returned (without accept/reject) as the SMCP3 weight; folding it into the particle weight gives (B),
thresholding $\log u<\mathcal W$ gives (A).

\paragraph{Marginal-likelihood estimator.} A genjax \code{ParticleCollection} reports
$\log\widehat Z = \lse_p(\log w^{(p)}) - \log P$ (\code{smc.py},
\code{get\_log\_marginal\_likelihood\_estimate}); the per-step product of these is the
$\log\widehat Z_W$ of Theorem~\ref{thm:smc}.

\begin{thebibliography}{9}
\bibitem{genjax} MIT Probabilistic Computing Project. \emph{genjax}: probabilistic programming with
programmable inference in JAX. \url{https://github.com/genjax-dev/genjax}.
\bibitem{gen} M.~F.~Cusumano-Towner, F.~A.~Saad, A.~K.~Lew, V.~K.~Mansinghka. Gen: a
general-purpose probabilistic programming system with programmable inference. \emph{PLDI}, 2019.
\bibitem{smcp3} A.~K.~Lew, G.~Matheos, et al. SMCP3: Sequential Monte Carlo with probabilistic
program proposals. \emph{AISTATS}, 2023.
\bibitem{chopin2020} N.~Chopin, O.~Papaspiliopoulos. \emph{An Introduction to Sequential Monte Carlo}.
Springer, 2020.
\end{thebibliography}

\end{document}
